% -- begin imports --
% import ./rotating_frame as rf
% import ./coupled_oscillators_hamiltonian as coh
% -- end imports --

\levelstay{Exercises}

\leveldown{Impulse response}
\begin{enumerate}[label=\alph*.]
  \item (Critical damping) Find the impulse response in the case $\omega_0 = \kappa / 2$.
  \item (Overdamping) Find the impulse response in the case $\omega_0 < \kappa / 2$.
  \item In this exercise, we prove statements from the text about the energy in the oscillator in the underdamped case.
    \begin{enumerate}[label=\roman*.]
      \item Write Newton's equation for a drivem damped mechanical harmonic oscillator with position $x$, mass $m$, spring constant $k$, coefficiency of friction $\mu$, and driving force $F(t)$.
        By comparison to the equations in the text, show that $\phi = x$, $\omega_0^2 = k/m$, $\kappa = \mu / m$, and $J(t) = F(t) / m$.
      \item Suppose the driving force is a delta function at time $t = t_0$, so the force is written $F(t) = \Delta P \, \delta(t - t_0)$ where $\Delta P$ is the impulse imparted to the system at time $t_0$.
        The driving term then has the form $J(t) = \Delta P \delta(t - t_0) / m$.
        By comparing with the main text where we write $J(t) = A \delta(t - t_0)$, interpret the meaning of $A$.
      \item Calculate the total energy
        \begin{equation*}
          E(t) = \frac{1}{2} k x(t)^2 + \frac{1}{2} m \dot{x}(t)^2
          \, .
        \end{equation*}
        Show that, in the limit $\kappa \ll \omega_0$, the energy is approximately
        \begin{equation*}
          E(t) = E_0 \exp \left( - \kappa t \right)
        \end{equation*}
        as stated in the main text.
        Show that $E_0 = m v_0^2 / 2$ where $v_0$ is the velocity immediately after time $t_0$.
    \end{enumerate}
\end{enumerate}


\levelstay{Solutions in the rotating frame}

\begin{enumerate}[label=\alph*.]
  \item Find the Green function for the driven oscillator described by Eq.~(\ref{rf.eq:equation_of_motion_phasor}).
  \item Solve for $\alpha(t)$ in the case that $J(t) = J_0 \delta(t)$. In this case, it's convenient to choose $\omega_f = \omega_0'$.
  \item Write down $\phi(t)$ for your solution to the previous part.
  \item Solve for $\alpha(t)$ in the case that $J(t) = J_0 \cos(\Omega t)\Theta(t)$ ($\Theta$ is the step function). Try both $\omega_f = \Omega$ and $\omega_f = \omega_0$.
  \item Write down $\phi(t)$ for your solution to the previous part.
\end{enumerate}

\levelstay{Coupled oscillators}

\begin{enumerate}[label=\alph*.]
  \item Derive Eqs.~(\ref{eq:coupled_oscillators_kirchhoff}) by application of Kirchhoff's laws.
  \item Check Eqs.~(\ref{coh.eq:coupled_oscillators_lagrangian}) by showing that Lagrange's equations of motion reproduce the equations of motion we got from Kirchhoff's laws.
\end{enumerate}
